\cvsection{Expériences professionnelles}

\begin{cventries}
  \cventry
    {Stagiaire ingénieur de recherche en IA}
    {INRIA}
    {Bordeaux, France}
    {Mars 2026 -- Septembre 2026}
    {
      \begin{cvitems}
        \item {Analyse du besoin des chercheurs et conception en autonomie d’une application fonctionnelle intégrant un agent IA multi-étapes pour la recherche d’articles scientifiques.}
        \item {Construction de l’agent sans framework d’orchestration avec API OpenAI, conception de prompts et gestion des outils, de l’état, de la mémoire et du contexte.}
        \item {Développement d’un backend Python/FastAPI avec API REST documentées avec OpenAPI, persistance MongoDB et interface React/Next.js.}
        \item {Mise en œuvre d’un pipeline RAG avec chunking, embeddings, recherche sémantique, pipeline d’ingestion et base vectorielle ChromaDB.}
        \item {Automatisation de la collecte de métadonnées scientifiques disponibles en ligne avec tâches cron et transmission des résultats par webhooks.}
        \item {Évaluation de LLM, notamment Qwen, Gemma et GLM, et de plusieurs quantifications selon la qualité, la latence, le TTFT et la taille du contexte.}
        \item {Utilisation de vLLM pour servir des modèles open source sur une infrastructure HPC équipée de GPU NVIDIA et gestion de modèles locaux avec LM Studio.}
        \item {Évaluation d’OpenClaw dans un environnement Docker isolé afin d’analyser les risques liés à l’accès d’un agent aux outils et aux fichiers.}
        \item {Conteneurisation avec Docker, gestion des sources avec Git et mise en place d’une CI/CD sur GitHub.}
        \item {Définition du protocole, analyse des résultats et rédaction en cours d’un article scientifique comparant les modèles et leurs limites.}
      \end{cvitems}
    }

  \cventry
    {Stagiaire ingénieur de recherche en apprentissage par renforcement}
    {Université Karunya}
    {Coimbatore, Inde}
    {Juin 2025 -- Septembre 2025}
    {
      \begin{cvitems}
        \item {Développement d’un modèle d’apprentissage par renforcement pour la navigation autonome d’un robot hospitalier, puis déploiement, optimisation et tests sur Raspberry Pi.}
      \end{cvitems}
    }
\end{cventries}
